\documentclass[ ../main.tex]{subfiles}
\providecommand{\mainx}{..}
\begin{document}
\section{Entropy maps}
\label{sec:entropy-maps}

We present a theoretical data type, denoted the \emph{entropy map}, that has several key properties:
\begin{enumerate}
\item Perfect model of the \emph{first-order} random approximate map.
\item Bit-length of in-place encoding is a random variable $\RV{L}$ that is expected to obtain the theoretical lower-bound.
\item In-place encoding has maximum conditional entropy (when conditioning on parameters like the expected false positive rate).
\end{enumerate}

\subsection{Framework}

The basic theory behind an entropy map is to map values in the domain to
values in the codomain by \emph{hashing} to a prefix-free code in the
codomain. We do not store anything related to the domain, since we are
simply hashing them, and a prefix of that hash will be used as a code
for a value in the codomain.

We consider sets in which the universe of elements is given by the \emph{countably infinite} set of all bit strings.
\begin{definition}
The countably infinite set of all bit strings is denoted by
\begin{equation}
    \BitSet^* = \left\{b \colon b \in \left\{0,1\right\}^*\right\}\,,
\end{equation}
where $*$ denotes all non-negative integers.
\end{definition}
A \emph{finite} subset of $\BitSet^*$ is given by the following definition.
\begin{definition}
The \emph{finite set} of all bit strings of length $n$ is denoted by
\begin{equation}
    \BitSet^n = \left\{b \colon b \in \left\{0,1\right\}^n\right\}
\end{equation}
with a cardinality given by
\begin{equation}
    \Card{\BitSet^n} = 2^n\,.
\end{equation}
\end{definition}

A \emph{hash function} is related to countable sets $\BitSet^*$ and $\BitSet^n$ and is given by the following definition.
\begin{definition}
A \emph{hash function} of type $\BitSet^* \mapsto \BitSet^n$ maps (hashes) bit strings of arbitrary-length to bit strings of fixed-length $n$.
\end{definition}

The random oracle is given by the following definition.
\begin{definition}
\label{def:randomoracle}
A random oracle $\hash \colon \BitSet^* \mapsto \BitSet^{\infty}$ is a theoretical \emph{hash function} whose output is uniformly distributed over its codomain.
\end{definition}

\begin{definition}[Entropy Map]
\label{def:entropy-map}
An \emph{entropy map} for a function $\Fun{f} \colon \Set{X} \mapsto \Set{Y}$
consists of:
\begin{itemize}
\item A hash function $\hash \colon \Set{X} \times \BitSet^{*} \mapsto \BitSet^{*}$.
\item A prefix-free code $\Fun{c} \colon \Set{Y} \mapsto \PS{\BitSet^{*}}$
      assigning sets of codewords to codomain values.
\item A decoding function $\Decode \colon \BitSet^{*} \mapsto \Maybe{\Set{Y}}$.
\end{itemize}
For each $y \in \Set{Y}$, the codewords in $\Fun{c}(y)$ form a prefix-free set.
\end{definition}

We allow many different codes for each value in the codomain, so that, for
instance, a code for the value $a$ may be $00$, $01$, $10$, and $11$.
We can efficiently decode this as $a$ if the hash is less than~4.

Suppose we have a prefix-free encoder-decoder for values of type $\Set{Y}$ denoted respectively by $\Encode_{\Set{Y}} \colon \Set{Y} \mapsto \BitSet^*$ and $\Decode_{\Set{Y}} \colon \BitSet^* \mapsto \Maybe{\Set{Y}}$.
If we are given a (possibly infinite) sequence of bits $b_1 b_2 \ldots b_k b_{k+1} \ldots$ and $b_1 b_2 \ldots b_k$ is a (prefix-free) code for some $y$ in $\Set{Y}$, then $\Decode_{\Set{Y}}(b_1 b_2 \ldots b_k b_{k+1} \ldots)$ maps to $y$.
Inversely, $\Encode_{\Set{Y}}(y)$ maps to $b_1 b_2 \ldots b_k$.

\subsection{Optimal space efficiency}

Suppose $\Prob{\Fun{f}(\RV{X}) = y} = p_y$, $\RV{X} \sim p_{\RV{X}}$, then the optimally space-efficient code, assuming a uniform hash function $\hash$, is to assign prefix-free codes for $y$ whose probability of being mapped to by $\hash$ sums to $p_y$.

\begin{theorem}[Optimal Code Length]
\label{thm:optimal-code}
For a uniform hash function $\hash$ and optimal prefix-free code assignment,
the expected bit length per element is
\begin{equation}
    \ell = H(\RV{Y}) = -\sum_{y \in \Set{Y}} p_y \log_2 p_y\,,
\end{equation}
where $H(\RV{Y})$ is the Shannon entropy~\cite{shannon1948} of the output distribution.
\end{theorem}
\begin{proof}
By Shannon's source coding theorem, the expected length
of an optimal prefix-free code for distribution $p_{\RV{Y}}$ equals the entropy
$H(\RV{Y})$. Since we assign codewords with total probability $p_y$ to each
$y \in \Set{Y}$, the result follows.
\end{proof}

If we imagine sampling $x$ from $\Set{X}$ with $p_{\RV{X}}$ and then mapping to
$y = \Fun{f}(x)$ and observing the sequence of $y$'s, then the expected bit length
is the entropy of the sequence of $y$'s. This is why we call it an
\emph{entropy map}.

For finite $\Set{X}$, we can imagine implicitly encoding the domain and
then for each value in the domain, storing the prefix-free code that it maps
to, which has an average bit length of $\ell$ and a total bit length of
$\Card{\Set{X}} \cdot \ell$.

\subsection{Rate-distortion}

The entropy map models a \emph{first-order} random approximate map where the approximation error is due to \emph{rate-distortion}~\cite{cover2006}.
Rate-distortion is a result of two different parameter choices:
\begin{enumerate}
	\item\label{itm:insufficient} Distortion due to \emph{insufficient} bits.
	\item\label{itm:fnr} Distortion due to accepting any false negative rate less than some lower-bound.
	\item Distortion due to specifying a false positive rate, which causes elements not in the set to be mapped to some uncertain value that is a function of the prefix-free encoder-decoder. If codes assume a uniform distribution, then randomly choose a $\RV{Y}$.
\end{enumerate}

In the singular hash map, \cref{itm:insufficient} is due to either a larger \emph{false positive rate} $\fprate = 2^{-k}$ by increasing $k$ or restricting the search to a maximum finite number of trials (and choosing the trial that resulted in the minimum loss according to some distance function).

In either case, the end result is an expected smaller encoding length, but in the first case we are directly motivated by a rate limit in the form of an upper-bound on the bit length and in the second case we are satisfied with greater approximation error but impose no upper-bound.

\subsection{Probability of success}

The set of rationals in $[0,1]$ is denoted by $\FancySet{P}$ and the set of power-of-two probabilities is denoted by
\begin{equation}
	\FancySet{R} \coloneqq \SetBuilder{2^{-k} \in \RatSet}{k \in \IntSet_{\geq 0}} \subset \FancySet{P}\,.
\end{equation}

The probability of success given a rate-distortion $\fnrate$ on positives and $\fprate$ on negatives is given by
\begin{equation}
	p(\fprate,\fnrate) = \binom{m}{(1-\fnrate)m} \fprate^{m(1-\fnrate)} (1-\fprate)^{m \fnrate}\,.
\end{equation}

Each trial is geometrically distributed, $\RV{T} \sim \geodist(p(\fprate,\fnrate))$.

The expected number of trials is $\Expect{\RV{T}} = \frac{1}{p(\fprate,\fnrate)}$.
The expected code size is $\Expect{\log_2 \RV{T}}$, which is non-linear and problematic but may be approximated with
\begin{equation}
	-\Fun{H}_b(\fnrate) - \fnrate \log_2 (1-\fprate) - (1-\fnrate) \log_2 \fprate
\end{equation}
bits per element where $\Fun{H}_b(p) \coloneqq (1-p)\log_2(1-p) + p \log_2 p$ is the binary entropy function.

This result is not the result of computing $\Expect{\log_2 \RV{T}}$, nor the expectation of its Taylor approximation.
It is simply $\log_2 \Expect{\RV{T}}$.

If we let $\fnrate \coloneqq 0$, the above simplifies to $\log_2 \fprate$, which is the lower-bound for \emph{first-order} random positive approximate sets over countably infinite universes with a false positive rate $\fprate$.

\subsection{Maybe monad}
Suppose the encoder-decoder has the pair $\Encode_{\Set{X}} \colon \Set{X} \mapsto \PS{\BitSet^*}$ and $\Decode_{\Set{X}} \colon \BitSet^* \pfun \Set{X}$.

We consider a specialization of the algorithm that, given a function $\Fun{f} \colon \Set{X} \mapsto \Set{Y}$ generates rate-distorted maps $\APFun{f}[\fprate][\fnrate] \colon \Set{X} \mapsto \Maybe{\Set{Y}}$.
The encoder-decoder only applies to elements of $\Set{X}$ which test as positive, at which point the elements are mapped to values as previously.

In other words, and assuming no rate-distortion, it maps values in the domain of definition of $\Fun{f}\!\!\restriction_{\Set{A}}$ to \emph{nothing} with probability $\fnrate$ and otherwise maps it correctly and it maps values not in its domain of definition to \emph{nothing} with probability $\fprate$ and value $y$ with probability
\begin{equation}
	p_{\RV{Y}}(y) \coloneqq \!\!\!\!\! \sum_{z \in \Encode_{\Set{Y}}(y)} \!\!\!\!\! 2^{-\BL(z)}\,.
\end{equation}

The function $\GenerateSeeds \colon (\Set{X} \mapsto \Maybe{\Set{Y})} \times \PS{\Set{X}} \times \FancySet{R} \mapsto \PS{\BitSet^* \times \FancySet{P}}$ is defined as
\begin{equation}
	\GenerateSeeds(\Fun{f}, \Set{A}, \fprate) \coloneqq \PlainSet{T}
	\end{equation}
where
\begin{equation}
\begin{split}
	m_a						&\coloneqq \Card{\Set{A}}\,,\\
	k						&\coloneqq -\log_2 \fprate\,,\\
	\Fun{h}(x \Given n,k)	&\coloneqq \trunc(\hash(x' \cat n'),k)\,,\\
	\Match(x \Given n,k)	&\coloneqq \Fun{h}(x \Given n,k)\,,\\
	\PlainSet{H} 			&\coloneqq \SetBuilder{\Match(x \Given n,k) \in \BitSet^{k}}{x \in \Set{X}}\,.
\end{split}
\end{equation}

\subsection{Construction algorithm}

\begin{algorithm}
	\caption{Implementation of \protect\MakeSingularHashMap over a universal set $\Set{X}$}
	\label{alg:makeset}
	\DontPrintSemicolon
	\SetKwProg{func}{function}{}{}
	\KwIn
	{
		$\Set{A}$ is a subset of a universal set $\Set{X}$.
		$\fprate$ is the false positive rate.
		$t$ is the maximum number of trials.
	}
	\KwOut
	{
		$\ASet{A}[\fprate]$, a rate-distorted approximate set of $\Set{A}$.
	}
	\func{\MakeSingularHashMap{$\Set{A}$, $\fprate$, $t$}}
	{
		$\Set{A}_{\BitSet^*} \coloneqq \SetBuilder{\Encode_{\Set{X}}(a) \in \BitSet^*}{a \in \Set{A}}$\;
		\For{$n \gets 0$ \KwTo $t$}
		{
			\For{$j \gets 1$ \KwTo $2^n$}
			{
				$\found \gets \True$\;
				$b_n \gets $ a bit string randomly drawn from $\BitSet^n$ without replacement\;
				$h_k \gets \Nothing$\;
				\For{$x \in \Set{A}_{\BitSet^*}$}
				{
					\If{$h_k = \Nothing$}
					{
						$h_k \gets \trunc(\hash(x \cat b_n), k)$\;
					}
					\ElseIf{$h_k \neq \trunc(\hash(x \cat b_n), k)$}
					{
						$\found \gets \False$\;
					}
				}
				\If{\found}
				{
					\Return $\Tuple{h_k, b_n}$\;
				}
			}
		}
	}
\end{algorithm}

\begin{figure}
    \centering
    %\input{img/fig_shmap.tex}
    \caption{\emph{Singular Hash Map}}
    \label{fig:shmap}
\end{figure}

The space required for the \emph{singular hash map} is of the order of the length $n$ of the bit string $b$. Therefore, for space efficiency, the algorithm exhaustively searches for a bit string in the order of increasing size $n$.

\subsection{Theoretical analysis}
The \emph{singular hash set} proves that the keys in $\APFun{f}$ are an \emph{approximate set} of the keys in $\Fun{f}$.

\begin{theorem}
The probability that a particular bit string $b \in \BitSet^*$ successfully codes the \emph{singular hash map} is given by
\begin{equation}
    p(m,\fprate,\mu) = \frac{\fprate^{m-1}}{2^{m \mu}}\,.
\end{equation}
\end{theorem}
\begin{proof}
For a particular bit string $b_n$, the joint probability that every key in $\Fun{f}$ collides and for each key $k$ in $\Fun{f}$, the concatenation of $1$ and $k$ has a hash in which the first $\BL(v)$ bits collide with mapped value $v$ is given by the following reasoning.

Suppose we have a set $\Fun{f} \coloneqq \{(k_1,v_1),\ldots,(k_m,v_m)\}$ and $k_1$ hashes to $y = \hash(k_1) \mod t$, where $\hash \colon \BitSet^* \mapsto \BitSet^N$ is a random hash function that uniformly distributes over its domain of $2^N$ possibilities and $N < \max{\BL(v_1),\ldots,\BL(v_m)}$.

The probability that $\hash(1 \cat k_1) \mod \BL(v_1) = v_1$ is given by
\begin{equation}
    \frac{1}{2^{\BL(v_1)}}\,.
\end{equation}
Since $y$ is a particular element in $\BitSet^k$, the probability that $k_j, j \neq 1$, hashes to $y$ is given by
\begin{equation}
    \frac{1}{2^k}\,.
\end{equation}
The probability that $1 \cat k_j, j \neq 1$ hashes to $v_j$ is given by
\begin{equation}
    \frac{1}{2^{\BL(v_j)}}\,.
\end{equation}
Since $\hash$ is a random hash function, every one of these probabilities are independent, and thus the joint probability is just the product given by
\begin{align}
    p &= \frac{1}{2^{k(m-1)}} \prod_{j=1}^{m} \frac{1}{2^{\BL(v_j)}}\\
      &= \fprate^{m-1} \frac{1}{2^{\sum_{j=1}^{m} \BL(v_j)}}\,.
\end{align}
The average bit length per key is given by
\begin{equation}
    \mu = \frac{1}{m} \sum_{j=1}^{m} \BL(v_j)
\end{equation}
and thus the probability $p$ may be parameterized as
\begin{equation}
    p = \frac{\fprate^{m-1}}{2^{m \mu}}\,.
\end{equation}
\end{proof}

\begin{theorem}
Given \emph{false positive} and \emph{false negative} domain error rates given by $\fprate$ and $\fnrate$ respectively, and a mean bit length $\mu$ of encodings of elements in the image of $\APFun{f}$, the \emph{expected} bit length of the Singular Hash Map asymptotically obtains a space complexity with a lower-bound given by
\begin{equation}
    -(1 - \fnrate) \log_2 \fprate + (1 - \fnrate) \mu \; \si{bits \per element}\,,
\end{equation}
which occurs when the encodings have uniformly distributed bit lengths, and an upper-bound given by
\begin{equation}
    -(1 - \fnrate) \log_2 \fprate + \mu \; \si{bits \per element}\,,
\end{equation}
which occurs when the encodings have degenerate bit lengths. The greater the entropy of the bit lengths, the smaller the expected bit length given a mean of $\mu$.
\end{theorem}
\begin{proof}
The space required for the singular hash set found by \Cref{alg:makeset} is of the order of the length $n$ of the bit string $b_n$.
Therefore, for space efficiency, the algorithm exhaustively searches for a bit string in the order of increasing size $n$.

We are interested in the first case when a success occurs, which is a geometric distribution with probability of success $p$ as given by
\begin{equation}
    \RV{Q} \sim \geodist\!\left(p = \frac{\fprate^{m-1}}{2^{m \mu}}\right)\,.
\end{equation}
The expected number of trials for the geometric distribution is given by
\begin{equation}
\label{eq:exp_trials}
    \Expect{\RV{Q}} = \frac{2^{m \mu}}{\fprate^{m-1}}\,.
\end{equation}
The $n$-th trial uniquely maps to a bit string of length $m = \lfloor \log_2 n \rfloor$.
Thus, the expected bit length is given approximately by
\begin{equation}
    \Expect{\log_2 \RV{Q}} = \log_2 \frac{2^{m \mu}}{\fprate^{m-1}}\; \si{bits}\,.
\end{equation}
Dividing by $m$ and simplifying results in
\begin{equation}
    -\frac{m-1}{m} \log_2 \fprate + \mu\; \si{bits \per element}\,.
\end{equation}
Asymptotically, as $m \to \infty$, the bits per element goes to
\begin{equation}
    -\log_2 \fprate + \mu\,.
\end{equation}
\end{proof}

\subsection{Trapdoor entropy maps}
\label{sec:trapdoor}

The smaller the bit lengths generated by the encoder, the smaller the entropy map.
However, optimal encoders \emph{reveal} information about the values in the entropy map.
For instance, if each input $x$ mapped to a list of numbers, then an optimal encoder will generate codes that minimize the length of the list, say a Huffman code.
However, by analyzing the Huffman codes produced by the encoder, we may infer the distribution of the lists.
Thus, for the sake of maintaining a \emph{trapdoor} state, the encoder should generate codes as though the values were \emph{uniformly} distributed.

A \emph{trapdoor} entropy map is an entropy map where the hash function is
applied to trapdoors (one-way transformations) of $\Set{X}$ and the
prefix-free codes for $\Set{Y}$ have no pattern, e.g., a
random selection of hash values for each value in $\Set{Y}$. This
provides a foundation for private set intersection, encrypted search, and
other secure computation primitives.

\subsection{Entropy}
\label{sec:entropy-map-entropy}

% TODO: expand this section with conditional entropy analysis
% H(f(x) | \tilde{f}(x)) quantifies remaining uncertainty after observation.

\end{document}
